\begin{explanationbox}
	\textbf{\textcolor{CExplanation}{为什么会想到这个结论？}}
	碰到 $2^{\log_2 8}$，先别急着算。$\log_2 8$ 的意思就是“$2$ 的几次方等于 $8$”？答案是 $3$，所以原式等于 $2^3=8$。
	\textbf{对数负责找指数，指数负责算出结果}，一来一回，数就还原了。

	\medskip
	\textbf{\textcolor{CExplanation}{用一张图看懂}}
	\begin{center}
	\begin{tikzpicture}[>=stealth, font=\small]
		\node[draw=CExplanation, rounded corners=3pt, minimum width=1.45cm, minimum height=0.82cm] (a) at (0,0) {$8$};
		\node[draw=CExplanation, rounded corners=3pt, minimum width=1.45cm, minimum height=0.82cm] (b) at (4.15,0) {$3$};
		\node[draw=CExplanation, rounded corners=3pt, minimum width=1.45cm, minimum height=0.82cm] (c) at (8.3,0) {$8$};
		\draw[->, thick, draw=CExplanation] (a) -- node[above, align=center] {找指数：$\log_2 8$} (b);
		\draw[->, thick, draw=CExplanation] (b) -- node[above, align=center] {算回去：$2^3$} (c);
		\node[font=\small] at (4.15,-0.95) {\textbf{先找指数，再乘方，回到原来的 $8$}};
	\end{tikzpicture}

	\end{center}

	\textbf{\textcolor{CExplanation}{反过来也一样}}
	先算 $2^3=8$，再问“$2$ 的几次方等于 $8$”，自然得到 $3$：
	\[
		\log_2(2^3)=\log_2 8=3.
	\]
	你可以把它想成：$5+3-3=5$；先做一步，再做相反的一步，就回到起点。

	\medskip
	\textbf{\textcolor{CExplanation}{什么时候用最快？}}
	\begin{itemize}[leftmargin=*, itemsep=0.35em, topsep=0.2em]
		\item \textbf{直接求值：}$2^{\log_2 3}=3$，不用求 $\log_2 3$ 的小数值。
		\item \textbf{先变形再求值：}$4^{\log_2 3}$ 不能直接还原，先把 $4$ 改写成 $2^2$。
		\item \textbf{解方程：}$5^{\log_5(2x-1)}=8$，先确认 $2x-1>0$，再化为 $2x-1=8$。
	\end{itemize}

	\textbf{\textcolor{CExplanation}{容易忽略的一点}}
	“还原”不是把式子里的 $\log$ 随便划掉。只有\textbf{两个底数一样}、\textbf{对数有意义}时，才可以这样做。
\end{explanationbox}
